\documentclass[12pt,a4paper]{article}

\usepackage[utf8]{inputenc}
\usepackage[T1]{fontenc}
\usepackage[italian]{babel}
\usepackage[margin=2.5cm]{geometry}
\usepackage{hyperref}
\usepackage{booktabs}
\usepackage{array}
\usepackage{enumitem}

\hypersetup{
    colorlinks=true,
    linkcolor=black,
    urlcolor=blue
}

\setlist[itemize]{topsep=4pt,itemsep=2pt,leftmargin=1.2cm}

\title{\textbf{Specifica di Progetto}\\[0.4cm]\large Study Planner}
\author{Samuele Caporale \\ 329096}
\date{Maggio 2026}

\begin{document}

\hypersetup{pageanchor=false}
\begin{titlepage}
    \centering
    {\Large Corso di Laurea in Informatica - scienza e tecnologia. \par}
    \vspace{1.2cm}
    {\huge \textbf{Specifica di Progetto}\par}
    \vspace{0.4cm}
    {\LARGE Study Planner\par}
    \vspace{1cm}
    {\large Applicazione web per l'organizzazione e la collaborazione sui progetti e task.\par}
    \vfill
    \begin{tabular}{rl}
        Studente: & Samuele Caporale \\
        Matricola: & 329096 \\
        Insegnamento: & Reti di calcolatori \\
        Data: & Maggio 2026 \\
    \end{tabular}
    \vfill
\end{titlepage}
\hypersetup{pageanchor=true}
\pagenumbering{arabic}

\section{Introduzione all'idea}

Il progetto \textit{Study Planner} nasce con l'obiettivo di realizzare una web app per la gestione dei task di studio in modo ordinato e collaborativo.
L'idea principale \`e di crare uno spazio in cui organizzare gruppi di studio/lavoro con milestone da raggiungere e task da svolgere in base alle priorità, mantenendo al tempo stesso una vista sintetica sulle metriche di processo.

Attualmente il progetto ha gi\`a una parte frontend avviata e funzionante, mentre l'evoluzione prevista riguarda soprattutto il consolidamento del backend, della persistenza cloud e della collaborazione real-time. La specifica proposta descrive quindi una soluzione gi\`a parzialmente sviluppata, ma pensata fin dall'inizio in modo modulare e scalabile.

\section{Tecnologie scelte}

\begin{table}[h]
    \centering
    \begin{tabular}{>{\raggedright\arraybackslash}p{3.2cm} >{\raggedright\arraybackslash}p{4.1cm} >{\raggedright\arraybackslash}p{6.1cm}}
        \toprule
        \textbf{Tecnologia} & \textbf{Ruolo} & \textbf{Motivazione sintetica} \\
        \midrule
        React 18 & Frontend SPA & Scelta per costruire un'interfaccia modulare e riusabile. \\
        Vite & Tooling frontend & Per la gestione del progetto e sviluppo. \\
        Zustand & State management & Per la gestione dello stato globale di supporto agli hooks di React. \\
        Bootstrap + Bootstrap Icons & UI e componenti & Per velocizzare la realizzazione della GUI seguendo il pattern "Responsive Web Design" e mantenendo una coerenza visiva. \\
        Chart.js & Analytics & Per la rappresentazione delle metriche di performance: statistiche e andamento dello studio/lavoro. \\
        Node.js & Backend API & Fornisce un ambiente unico JavaScript lato server, coerente con il frontend. \\
        localStorage & Persistenza attuale & Utile in fase di prototipazione per testare il frontend senza bisogno di un backend sin dall'inizio. \\
        Firebase Storage & Storage cloud  & Permette di estendere il progetto verso la persistenza remota, accesso multi-dispositivo. \\
        WebSocket / Socket.IO & Real-time & Rende possibile la sincronizzazione immediata dei task condivisi tra più utenti. \\
        \bottomrule
    \end{tabular}
\end{table}

\subsection*{Agenti AI}
Un aspetto rilevante del progetto \`e l'uso di agenti AI come supporto allo sviluppo. L'approccio adottato non \`e quello di delegare completamente la realizzazione tecnica, ma di lavorare in modo pi\`u vicino al ruolo di \textit{project manager} e \textit{software architect}: si definiscono struttura, vincoli, priorit\`a e obiettivi, mentre gli agenti aiutano nella produzione rapida di componenti e refactoring.
Le parti generate tramite questo approccio vengono comunque riviste, corrette e rifinite manualmente, soprattutto nei dettagli di interfaccia, nella coerenza architetturale e nelle scelte finali di integrazione. In questo senso il vibecoding \`e stato usato come acceleratore di produttivit\`a.

Per fare ciò sono stati creati agenti specifici per:
\begin{itemize}
    \item Progettazione UX/UI del frontend.
    \item Implementazione di componenti React e logica di stato.
    \item Cordinatore e validatore dell'operato dei due precedenti, e testing.
\end{itemize}

Ognuno dei quali è tenuto a rispettare pattern di progettazione, e le migliori pratiche di sviluppo s/w attraverso l'utilizzo di Skills.
\begin{itemize}
    \item L'agente di progettazione UX/UI si occupa di definire l'aspetto visivo e l'esperienza utente, creando mockup e prototipi.
    \item L'agente di implementazione si concentra sulla scrittura del codice, seguendo le specifiche fornite dall'agente di progettazione e rispettando le linee guida architetturali.
    \item L'agente coordinatore monitora il lavoro degli altri due, assicurandosi che le componenti siano integrate correttamente.
\end{itemize}

\section{Architettura}

L'architettura scelta segue una separazione netta tra frontend e backend. Il frontend si occupa della presentazione, dell'interazione utente e della gestione dello stato client, mentre il backend \`e pensato come livello dedicato a API, persistenza e servizi di collaborazione.

\section{Frontend}

Il frontend rappresenta attualmente la parte pi\`u avanzata del progetto. Si tratta di una Single Page Application sviluppata in React, organizzata in componenti, hook, servizi e store globale.

Tra gli elementi gi\`a rilevanti del frontend si possono evidenziare:
\begin{itemize}
    \item dashboard centrale per la gestione dei task;
    \item organizzazione dei task in gruppi e milestone;
    \item timer di focus integrato per il lavoro su singole attivit\`a;
    \item sezione analytics con grafici e statistiche;
    \item interfaccia responsive con gestione tema e componenti riusabili.
\end{itemize}

Dal punto di vista progettuale, il frontend \`e stato sviluppato per validare rapidamente l'esperienza utente e definire una base concreta su cui innestare le funzionalit\`a distribuite.

\section{Backend}

Il backend \`e impostato come servizio separato in ambiente Node.js, con una struttura pensata per ospitare API REST, moduli funzionali e integrazioni successive. Anche se questa parte \`e meno avanzata del frontend, la sua presenza architetturale \`e fondamentale per superare i limiti di una soluzione solo client-side.

I punti chiave previsti per il backend sono:
\begin{itemize}
    \item esposizione di endpoint REST per task, gruppi e utenti;
    \item gestione centralizzata della logica applicativa;
    \item supporto a autenticazione e autorizzazione;
    \item integrazione con storage persistente e servizi real-time.
\end{itemize}

\section{Cloud Storage}
Per la fase iniziale il progetto utilizza una persistenza locale, utile per prototipazione e test veloci. Tuttavia, per rendere l'applicazione realmente distribuita, \`e prevista l'introduzione di uno strato cloud dedicato allo storage.
L'uso di uno storage cloud permette di:
\begin{itemize}
    \item conservare i dati in modo persistente.
    \item abilitare l'accesso a più utenti.
    \item predisporre meccanismi di backup ed eventuale esportazione dati.
    \item supportare un'evoluzione del progetto verso scenari collaborativi.
\end{itemize}

La scelta di un servizio gestito come Firebase Storage, \`e motivata dalla rapidit\`a di integrazione e dalla riduzione della complessit\`a infrastrutturale.

\section{Socket e collaborazione real-time}
Una delle evoluzioni pi\`u interessanti del progetto riguarda la comunicazione real-time tramite socket.
L'obiettivo \`e permettere la condivisione dei task e l'aggiornamento immediato delle modifiche tra utenti collegati allo stesso progetto.
Questa componente si occuperà di:
\begin{itemize}
    \item sincronizzare la creazione o modifica dei task in tempo reale.
    \item mostrare aggiornamenti immediati sul completamento delle attività.
    \item coordinare funzionalit\`a come timer o avanzamento del lavoro comune.
\end{itemize}

\section{Conclusione}
In sintesi, il progetto da realizzare rappresenta una Web App per la pianificazione, condivisione ed analisi di piccoli task giornalieri che se completati vanno ad avanzare lo stato delle milestone fino al raggiungimento dell'obiettivo finale. 
Tutto necessita della persistenza dei dati in cloud e collaborazione real-time. 
L'uso di agenti AI si inserisce in questo contesto come supporto operativo, mantenendo comunque centrale il ruolo umano nelle decisioni progettuali e nelle rifiniture finali.
\end{document}
